\chapter{Twierdzenie o funkcji uwikłanej}
\label{ma2:ch:implicit-function}

Niech $X,Y,Z$ będą rzeczywistymi przestrzeniami Banacha.
Na $X\times Y$ przyjmujemy normę maksimum
$\|(x,y)\|=\max\{\|x\|,\|y\|\}$, względem której
produkt jest przestrzenią Banacha: ciąg Cauchy'ego ma
obie składowe będące ciągami Cauchy'ego, a ich granice
wyznaczają granicę w normie maksimum. Niech
$\mathcal U\subseteq X\times Y$ będzie zbiorem otwartym.
Równanie $F(x,y)=0$ może wyznaczać $y$ jako funkcję $x$,
chociaż nie znamy jawnego wzoru rozwiązania.

\section{Istnienie i jednoznaczność rozwiązania lokalnego}

Dla $F:\mathcal U\to Z$ różniczkowalnego w $(x,y)$ oznaczamy
\[
 D_xF(x,y)h=DF(x,y)[(h,0)],\qquad
 D_yF(x,y)k=DF(x,y)[(0,k)].
\]
Są to ograniczone operatory odpowiednio z $X$ do $Z$
i z $Y$ do $Z$: włożenia $h\mapsto(h,0)$ oraz
$k\mapsto(0,k)$ są liniowe i ograniczone.
Z liniowości $DF(x,y)$ wynika
\[
 DF(x,y)[(h,k)]=D_xF(x,y)h+D_yF(x,y)k.
\]
Jeżeli $F$ jest klasy $C^1$, oba operatory częściowe
zależą w sposób ciągły od $(x,y)$ w normie operatorowej,
ponieważ dla dwóch operatorów $L,M\in\mathcal L(X\times Y,Z)$
oraz włożenia $i_X:h\mapsto(h,0)$ mamy
$\|(L-M)i_X\|\le\|L-M\|$.
Analogicznie dla włożenia $i_Y:k\mapsto(0,k)$.
Stosujemy te oszacowania do wartości $DF$ w dwóch punktach.

\begin{twierdzenie}[Twierdzenie o funkcji uwikłanej]
\label{ma2:thm:implicit-function}
Niech $F\in C^1(\mathcal U,Z)$, $(a,b)\in\mathcal U$,
$F(a,b)=0$ i niech $A=D_yF(a,b):Y\to Z$ będzie
izomorfizmem liniowym i topologicznym.
Istnieją otwarte otoczenia $U_0\subseteq X$ punktu $a$
i $V_0\subseteq Y$ punktu $b$ takie, że
$U_0\times V_0\subseteq\mathcal U$, oraz jedyna funkcja
$\varphi:U_0\to V_0$ spełniająca
\[
 F(x,y)=0\quad\Longleftrightarrow\quad y=\varphi(x)
 \qquad ((x,y)\in U_0\times V_0).
\]
Mamy $\varphi(a)=b$, $\varphi\in C^1(U_0,Y)$,
a $D_yF(x,\varphi(x))$ jest izomorfizmem dla każdego $x\in U_0$.
Ponadto
\[
 D\varphi(x)=-(D_yF(x,\varphi(x)))^{-1}
                  D_xF(x,\varphi(x)).
\]
Jeżeli $F$ jest klasy $C^m$, gdzie $m\ge2$,
lub klasy $C^\infty$, to $\varphi$ jest tej samej klasy.
\end{twierdzenie}

Dowód twierdzenia znajduje się w
\atlasref{FA-15::thm:implicit15}{FA-15: twierdzenie o funkcji uwikłanej}.
Jego hipotezy są dokładnie hipotezami przyjętymi tutaj.
Przedstawmy konstrukcję wiążącą go z poprzednim rozdziałem.
Definiujemy $\Psi:\mathcal U\to X\times Z$ przez
$\Psi(x,y)=(x,F(x,y))$. W punkcie $(a,b)$, przy
$B=D_xF(a,b)$, otrzymujemy
\[
 D\Psi(a,b)[(h,k)]=(h,Bh+Ak).
\]
Odwrotnością tego operatora jest
operator $L:X\times Z\to X\times Y$ dany wzorem
$L(h,z)=(h,A^{-1}(z-Bh))$. Istotnie,
\[
 \begin{aligned}
 L(D\Psi(a,b)[(h,k)])&=(h,A^{-1}Ak)=(h,k),\\
 D\Psi(a,b)[L(h,z)]&=(h,Bh+z-Bh)=(h,z).
 \end{aligned}
\]
Operator $L$ jest liniowy i ograniczony, gdyż
\[
 \|A^{-1}(z-Bh)\|
 \le\|A^{-1}\|(\|z\|+\|B\|\|h\|).
\]
Stąd $\|L(h,z)\|\le
\max\{1,\|A^{-1}\|(1+\|B\|)\}\max\{\|h\|,\|z\|\}$.
Oba produkty są przestrzeniami Banacha, a $\Psi$
jest klasy $C^1$ (pierwsza składowa jest rzutem).
Twierdzenie \ref{ma2:thm:inverse-function} stosuje się
więc do $\Psi$. Odwrotność lokalna w otoczeniu $(a,0)$
ma postać $(x,z)\mapsto(x,G(x,z))$, ponieważ pierwsza
współrzędna $\Psi$ jest niezmieniona.
Wybierając otoczenia produktowe i zmniejszając otoczenie
$a$ tak, by $G(x,0)$ pozostawało w otoczeniu $b$,
otrzymujemy $\varphi(x)=G(x,0)$. Szczegóły doboru
otoczeń i jednoznaczności podano w przywołanym dowodzie.
Ponieważ $F(a,b)=0$ i $b\in V_0$, równoważność
z tezy daje $\varphi(a)=b$. Dla każdego $x\in U_0$
różniczkujemy tożsamość $F(x,\varphi(x))=0$.
Z reguły łańcuchowej (twierdzenie \ref{ma2:thm:frechet-chain})
dla $h\in X$ otrzymujemy
\[
 0=DF(x,\varphi(x))[(h,D\varphi(x)h)]
  =D_xF(x,\varphi(x))h+
    D_yF(x,\varphi(x))D\varphi(x)h.
\]
Mnożenie przez odwracalny operator
$(D_yF(x,\varphi(x)))^{-1}:Z\to Y$
daje wzór na $D\varphi(x)$.
Odwracalność pochodnej częściowej wzdłuż wykresu
zapewnia przywołany dowód FA-15; wynika ona
z ciągłości pochodnej i oszacowania perturbacyjnego
po ewentualnym zmniejszeniu $U_0$.

\begin{stwierdzenie}[Równania z parametrami i prawą stroną]
\label{ma2:prop:implicit-parameter}
Przy hipotezach twierdzenia \ref{ma2:thm:implicit-function}
istnieją otwarte otoczenia $U_1$ punktu $a$, $W_1$ zera w $Z$
i $V_1$ punktu $b$ oraz funkcja
$G:U_1\times W_1\to V_1$ klasy $C^1$ takie, że
$U_1\times V_1\subseteq\mathcal U$ i
\[
 F(x,y)=z\quad\Longleftrightarrow\quad y=G(x,z)
 \quad(x\in U_1,\ z\in W_1,\ y\in V_1).
\]
Dla $(x,z)\in U_1\times W_1$ oznaczmy $q=(x,G(x,z))$.
Operator $D_yF(q)$ jest izomorfizmem i mamy
\[
 D_xG(x,z)=-(D_yF(q))^{-1}D_xF(q),\qquad
 D_zG(x,z)=(D_yF(q))^{-1}.
\]
Jeżeli $F$ jest klasy $C^m$ ($m\ge2$) lub $C^\infty$,
to $G$ jest tej samej klasy.
\end{stwierdzenie}
\begin{proof}
Na otwartym zbiorze
$\{((x,z),y):(x,y)\in\mathcal U,\ z\in Z\}$
w $(X\times Z)\times Y$ definiujemy
$H((x,z),y)=F(x,y)-z$. Jest to zbiór otwarty jako
przeciwobraz $\mathcal U$ przez ciągły rzut na $(x,y)$.
Funkcja $H$ jest klasy $C^1$ przez liniowość rzutów i
regułę łańcuchową. Mamy $H((a,0),b)=0$ oraz
$D_yH((a,0),b)=A$. Wszystkie hipotezy twierdzenia
\ref{ma2:thm:implicit-function} są spełnione.
Otrzymujemy funkcję $G$ na otoczeniu $(a,0)$,
które zawiera otoczenie produktowe $U_1\times W_1$.
Z twierdzenia otrzymujemy również otwarte otoczenie
$V_1$ punktu $b$. Ponieważ $0\in W_1$, dla
$(x,y)\in U_1\times V_1$ punkt $((x,0),y)$ należy
do dziedziny $H$, skąd $(x,y)\in\mathcal U$.
Dla wszystkich $(x,z,y)\in U_1\times W_1\times V_1$
równoważność z twierdzenia daje
$F(x,y)=z$ wtedy i tylko wtedy, gdy $y=G(x,z)$,
bo $H((x,z),y)=F(x,y)-z$.
Wzór na pochodną funkcji uwikłanej zastosowany do $H$
daje
\[
 DG(x,z)[(h,\ell)]=-(D_yF(q))^{-1}(D_xF(q)h-\ell).
\]
Odwracalność $D_yF(q)=D_yH((x,z),G(x,z))$
wynika z tezy tego twierdzenia. Podstawiamy kolejno
$\ell=0$ i $h=0$, otrzymując oba żądane operatory
częściowe. Regularność
wyższych rzędów wynika z odpowiedniej regularności $H$
oraz tego samego twierdzenia.
\end{proof}

\section{Pochodne funkcji uwikłanej}

\begin{stwierdzenie}[Druga pochodna]
\label{ma2:prop:implicit-second}
Przy hipotezach twierdzenia \ref{ma2:thm:implicit-function}
załóżmy dodatkowo, że $F\in C^2(\mathcal U,Z)$.
Dla $x\in U_0$ oznaczmy $q=(x,\varphi(x))$,
$A_x=D_yF(q)$ i $P_x=D\varphi(x)$.
Wtedy dla $h,k\in X$
\[
 D^2\varphi(x)[h,k]
 =-A_x^{-1}D^2F(q)[(h,P_xh),(k,P_xk)].
\]
Przy normie maksimum na produkcie zachodzi
\[
 \|D^2\varphi(x)\|
 \le\|A_x^{-1}\|\|D^2F(q)\|
                  \max\{1,\|P_x\|\}^2.
\]
\end{stwierdzenie}
\begin{proof}
Funkcja $\varphi$ jest klasy $C^2$ na mocy twierdzenia.
Niech $\gamma(x)=(x,\varphi(x))$. Z liniowości
różniczkowania składowych mamy
$D\gamma(x)h=(h,P_xh)$ oraz
$D^2\gamma(x)[h,k]=(0,D^2\varphi(x)[h,k])$.
Zachodzi $F\circ\gamma=0$, a $\gamma(U_0)\subseteq\mathcal U$.
Stosujemy wzór na drugą pochodną złożenia
(stwierdzenie \ref{ma2:prop:second-chain}) do tych
dwóch odwzorowań klasy $C^2$. Otrzymujemy
\[
 \begin{aligned}
 0={}&D^2F(q)[(h,P_xh),(k,P_xk)]\\
     &+DF(q)[(0,D^2\varphi(x)[h,k])]\\
   ={}&D^2F(q)[(h,P_xh),(k,P_xk)]
      +A_xD^2\varphi(x)[h,k].
 \end{aligned}
\]
Operator $A_x$ jest odwracalny, więc działamy na
równość przez $A_x^{-1}$ i przenosimy pierwszy składnik.
W normie maksimum
$\|(h,P_xh)\|\le\max\{1,\|P_x\|\}\|h\|$,
i analogiczne oszacowanie zachodzi dla $k$.
W konsekwencji
\[
 \begin{aligned}
 \|D^2\varphi(x)[h,k]\|
 &\le\|A_x^{-1}\|\|D^2F(q)\|
       \|(h,P_xh)\|\|(k,P_xk)\|\\
 &\le\|A_x^{-1}\|\|D^2F(q)\|
       \max\{1,\|P_x\|\}^2\|h\|\|k\|.
 \end{aligned}
\]
Biorąc supremum po $\|h\|,\|k\|\le1$, otrzymujemy
oszacowanie normy dwuliniowej.
\end{proof}

\begin{wniosek}[Jedna zmienna zależna i jedna niezależna]
\label{ma2:cor:implicit-scalar}
Dla $F:\mathcal U\subseteq\R^2\to\R$ klasy $C^2$
i rozwiązania lokalnego $y=\varphi(x)$ z twierdzenia
\ref{ma2:thm:implicit-function} zachodzi
\[
 \begin{aligned}
 \varphi'(x)&=-\frac{F_x}{F_y},\\
 \varphi''(x)&=-\frac{F_{xx}+2F_{xy}\varphi'(x)
                          +F_{yy}(\varphi'(x))^2}{F_y},
 \end{aligned}
\]
gdzie wszystkie pochodne $F$ obliczamy w $(x,\varphi(x))$.
\end{wniosek}
\begin{proof}
Operator $D_yF$ jest mnożeniem przez niezerową liczbę
$F_y$. Pierwszy wzór wynika więc bezpośrednio z twierdzenia.
W drugim wzorze podstawiamy $h=k=1$ w stwierdzeniu
\ref{ma2:prop:implicit-second}. Macierz drugiej
pochodnej $F$ jest macierzą Hessego.
Ponieważ $F$ jest klasy $C^2$, mieszane pochodne
są równe na mocy twierdzenia \ref{ma2:thm:schwarz}.
Wartość formy dwuliniowej na parze
$(1,\varphi'(x)),(1,\varphi'(x))$ wynosi
$F_{xx}+2F_{xy}\varphi'(x)+F_{yy}(\varphi'(x))^2$.
\end{proof}

\section{Interpretacja geometryczna}

\begin{przyklad}[Sfera jako lokalny wykres]
\label{ma2:ex:implicit-sphere}
Rozważmy $F(x_1,x_2,y)=x_1^2+x_2^2+y^2-1$.
W punkcie $(0,0,1)$ mamy $F=0$ i $D_yF=2$, więc
twierdzenie \ref{ma2:thm:implicit-function} daje lokalną
funkcję klasy $C^\infty$. Na otwartym dysku
$D=\{x\in\R^2:\|x\|_2<1\}$ rozwiązanie dodatnie jest
jednoznaczne i ma postać
$\varphi(x)=\sqrt{1-\|x\|_2^2}$, ponieważ
$y^2=1-\|x\|_2^2>0$.
Funkcja $\varphi$ jest gładka na $D$, ponieważ
pierwiastek jest gładki na $(0,\infty)$, a jego
argument jest dodatnim wielomianem na $D$.
W każdym punkcie $(x,\varphi(x))$ operator $D_yF$
jest mnożeniem przez $2\varphi(x)>0$; wzór z twierdzenia
można więc zastosować lokalnie w każdym takim punkcie,
a jednoznaczność rozwiązania dodatniego utożsamia
lokalne rozwiązanie z ograniczeniem $\varphi$.
Dla $h\in\R^2$ mamy $D_xF(x,y)h=2\langle x,h\rangle$,
a $D_yF(x,y)k=2yk$. Zatem
\[
 D\varphi(x)h=-\frac{\langle x,h\rangle}{\varphi(x)}.
\]
W szczególności $D\varphi(0)=0$. Parametryzacja
$x\mapsto(x,\varphi(x))$ ma w zerze pochodną
$h\mapsto(h,0)$ o obrazie $\R^2\times\{0\}$.
Przesunięcie tego obrazu przez punkt $(0,0,1)$
daje płaszczyznę styczną $y=1$.
Różniczkując $DF(x,y)[(k,t)]=2\langle x,k\rangle+2yt$
w kierunku $(h,s)$, widzimy, że druga pochodna $F$
jest stałą formą
$D^2F(x,y)[(h,s),(k,t)]=2(\langle h,k\rangle+st)$.
Ze stwierdzenia \ref{ma2:prop:implicit-second} wynika
\[
 \begin{aligned}
 D^2\varphi(x)[h,k]
 &=-\frac{\langle h,k\rangle+
                   D\varphi(x)h\,D\varphi(x)k}{\varphi(x)}\\
 &=-\frac{\langle h,k\rangle}{\varphi(x)}
 -\frac{\langle x,h\rangle\langle x,k\rangle}{\varphi(x)^3}.
 \end{aligned}
\]
W szczególności w punkcie $x=0$ mamy $D^2\varphi(0)[h,k]=-\langle h,k\rangle$.
Ze wzoru Taylora (twierdzenie \ref{ma2:thm:taylor-peano})
otrzymujemy $\varphi(x)=1-\tfrac12\|x\|_2^2+o(\|x\|_2^2)$
przy $x\to0$.
Na równiku, gdzie $y=0$, nie można zastosować
twierdzenia do wyznaczania tej współrzędnej: $D_yF=0$.
Co więcej, w żadnym otoczeniu punktu równika
$(x^0,0)$, gdzie $\|x^0\|_2=1$, sfera nie jest
wykresem jednej funkcji $y$ zmiennych $x_1,x_2$.
Dla $0<\varepsilon<1$ punkty
\[
 ((1-\varepsilon)x^0,\ \pm\sqrt{2\varepsilon-\varepsilon^2})
\]
leżą na sferze, mają tę samą składową $x$ i różne
składowe $y$. Oba dążą do $(x^0,0)$ przy
$\varepsilon\to0^+$, więc należą do każdego jego
otwartego otoczenia dla dostatecznie małego
$\varepsilon$. Jeżeli jednak $x_1^0\ne0$, to
$D_{x_1}F(x^0,0)=2x_1^0\ne0$ i lokalnie można
wyznaczyć $x_1$ jako funkcję $(x_2,y)$.
Jeżeli $x_1^0=0$, to $|x_2^0|=1$ i analogicznie
wyznaczamy $x_2$ jako funkcję $(x_1,y)$.
\end{przyklad}

\paragraph{Rysunek 3D: sfera, lokalny wykres i płaszczyzna styczna.}
Na sferze jednostkowej wyróżniamy fragment górnej
półsfery nad dyskiem o promieniu $0.65$.
Pozostałą część sfery rysujemy dla kąta biegunowego
$\alpha\in[\arcsin(0.65),\pi]$, aby nie nakładać
dwóch powierzchni na wyróżnionym fragmencie. Płaszczyzna
$y=1$ jest pokazana nad tym samym dyskiem; nie jest
częścią sfery poza punktem styczności. Kod tworzy
interaktywny wykres Plotly. W kodzie oś pionową
nazywamy $y$, a dwie poziome $x_1,x_2$.
Rysunek \ref{ma2:fig:implicit-sphere} przedstawia tę samą
konstrukcję w wersji statycznej.

\begin{figure}[htbp]
\centering
\begin{tikzpicture}[x={(2.3cm,-0.65cm)},y={(1.4cm,0.95cm)},z={(0cm,2.4cm)},>=Stealth,font=\small]
\pgfmathsetmacro{\capangle}{asin(0.65)}
\foreach \a in {0,30,...,330}{
  \draw[gray!65,thin] plot[domain=\capangle:180,samples=42,variable=\t]
    ({sin(\t)*cos(\a)},{sin(\t)*sin(\a)},{cos(\t)});
}
\foreach \a in {45,60,90,120,150}{
  \draw[gray!65,thin] plot[domain=0:360,samples=65,variable=\t]
    ({sin(\a)*cos(\t)},{sin(\a)*sin(\t)},{cos(\a)});
}
\fill[orange!50,fill opacity=.28] plot[domain=0:360,samples=73,variable=\t]
  ({.65*cos(\t)},{.65*sin(\t)},1) -- cycle;
\draw[orange!80!black,thick] plot[domain=0:360,samples=73,variable=\t]
  ({.65*cos(\t)},{.65*sin(\t)},1);
\foreach \a in {0,30,...,330}{
  \draw[blue!70!black,thin] plot[domain=0:\capangle,samples=23,variable=\t]
    ({sin(\t)*cos(\a)},{sin(\t)*sin(\a)},{cos(\t)});
}
\foreach \r in {.22,.44,.65}{
  \draw[blue!70!black,thick] plot[domain=0:360,samples=73,variable=\t]
    ({\r*cos(\t)},{\r*sin(\t)},{sqrt(1-\r*\r)});
}
\draw[->,gray!80!black] (0,0,0)--(1.45,0,0) node[right] {$x_1$};
\draw[->,gray!80!black] (0,0,0)--(0,1.45,0) node[right] {$x_2$};
\draw[->,gray!80!black] (0,0,-1.15)--(0,0,1.4) node[above] {$y$};
\fill[black] (0,0,1) circle[radius=1.5pt];
\draw[orange!80!black,thin] (.65,0,1)--(1.15,0,1.25) node[right] {$y=1$};
\draw[blue!70!black,thin] (0,.65,{sqrt(1-.65*.65)})--(0,1.5,.75) node[right] {$y=\varphi(x)$};
\end{tikzpicture}
\caption{Sfera jednostkowa: lokalny wykres (niebieski), jego płaszczyzna styczna w $(0,0,1)$ (pomarańczowa) i pozostała część sfery (szara).}
\label{ma2:fig:implicit-sphere}
\end{figure}

\begin{lstlisting}[language=Python]
import numpy as np
import plotly.graph_objects as go

a = np.linspace(np.arcsin(0.65), np.pi, 61)
t = np.linspace(0, 2*np.pi, 121)
A, T = np.meshgrid(a, t)
fig = go.Figure()
fig.add_trace(go.Surface(
    x=np.sin(A)*np.cos(T), y=np.sin(A)*np.sin(T),
    z=np.cos(A), opacity=0.25, colorscale="Greys",
    showscale=False, name="sphere remainder"))
r = np.linspace(0, 0.65, 41)
R, T = np.meshgrid(r, t)
X, Y = R*np.cos(T), R*np.sin(T)
fig.add_trace(go.Surface(x=X, y=Y,
    z=np.sqrt(1-R**2), colorscale="Blues",
    showscale=False, name="local graph"))
fig.add_trace(go.Surface(x=X, y=Y,
    z=np.ones_like(R), opacity=0.45,
    colorscale="Oranges", showscale=False,
    name="tangent plane"))
fig.update_layout(scene=dict(aspectmode="data",
    xaxis_title="x1", yaxis_title="x2", zaxis_title="y"))
fig.write_html("am2_19_implicit.html", include_plotlyjs=True)
fig.show()
\end{lstlisting}

\begin{przyklad}[Brak jednoznaczności i brak istnienia]
\label{ma2:ex:implicit-degenerate}
Niech $F(x,y)=y^2-x$. W punkcie $(0,0)$ mamy
$F(0,0)=0$, lecz $D_yF(0,0)=0$. Każde otwarte
otoczenie $U_0$ zera w $\R$ zawiera pewien przedział
$(-\delta,\delta)$, a każde otwarte otoczenie $V_0$
zera zawiera $(-\varepsilon,\varepsilon)$.
Dla $0<x<\min\{\delta,\varepsilon^2\}$ istnieją dwa
różne rozwiązania $y=\pm\sqrt{x}$ w $V_0$,
natomiast dla $-\delta<x<0$ nie istnieje żadne
rozwiązanie rzeczywiste. Nie ma więc funkcji
spełniającej równoważność z twierdzenia w otoczeniu zera.
Z kolei dla $F(x,y)=(y-x)^3$ mamy
$F(x,y)=0$ wtedy i tylko wtedy, gdy $y-x=0$.
Zbiór zer jest więc globalnie wykresem gładkiej
funkcji $\varphi(x)=x$, chociaż
$D_yF(x,y)=3(y-x)^2$, a zatem $D_yF(x,x)=0$
dla wszystkich $x$.
Warunek odwracalności pochodnej częściowej jest
dostateczny, lecz nie jest konieczny do istnienia
gładkiego wykresu zbioru zer.
\end{przyklad}

\section{Przykład operatorowy}

\begin{przyklad}[Równanie zależne od parametru]
\label{ma2:ex:implicit-operator}
Niech $E$ będzie rzeczywistą przestrzenią Banacha,
$K\in\mathcal L(E)$ i $b\in E$. Rozważmy
$F(t,u)=u-tKu-b$ na $\R\times E$.
Przyrost $F$ w kierunku $(s,h)$ jest równy
$(I-tK)h-sKu-sKh$. Ostatni składnik ma normę
nie większą niż $\|K\|\max\{|s|,\|h\|\}^2$,
a po podzieleniu przez normę przyrostu niezerowego
otrzymujemy oszacowanie dążące do zera.
Część liniowa jest ograniczona, ponieważ
\[
 \|(I-tK)h-sKu\|\le
 (1+|t|\|K\|+\|K\|\|u\|)\max\{|s|,\|h\|\}.
\]
Zatem $DF(t,u)[(s,h)]=(I-tK)h-sKu$.
Przyrost $DF$ w kierunku $(r,k)$ jest operatorem
$(s,h)\mapsto-rKh-sKk$ o normie co najwyżej
$2\|K\|\max\{|r|,\|k\|\}$. Zależność od $(r,k)$
jest liniowa, a reszta jest zerowa. Dlatego $DF$
jest ciągłe, a jego pochodna jest stała:
$D^2F(t,u)[(s,h),(r,k)]=-sKk-rKh$,
a wyższe pochodne są zerowe. Zatem $F$ jest klasy
$C^\infty$.
Mamy $F(0,b)=0$ i $D_uF(0,b)=I_E$, więc istnieje
lokalne rozwiązanie $u(t)$ klasy $C^\infty$.
Z równania $(I-tK)u(t)=b$ przez różniczkowanie
(reguła dla ciągłego odwzorowania dwuliniowego
$(T,w)\mapsto Tw$, dla którego $\|Tw\|\le\|T\|\|w\|$;
twierdzenie \ref{ma2:thm:frechet-bilinear})
otrzymujemy najpierw $-Ku(t)+(I-tK)u'(t)=0$.
Różniczkując ponownie, dostajemy
$-Ku'(t)+(I-tK)u''(t)=Ku'(t)$. Stąd
\[
 (I-tK)u'(t)=Ku(t),\qquad
 (I-tK)u''(t)=2Ku'(t).
\]
Ponieważ $u(0)=b$, w punkcie $t=0$ pierwsza
równość daje $u'(0)=Kb$, a druga
$u''(0)=2Ku'(0)=2K^2b$.
Jeśli $E=\{0\}$, równanie ma jedyne rozwiązanie
zerowe dla każdego $t\in\R$.
Jeśli $E\ne\{0\}$ oraz $|t|\|K\|<1$, stosujemy
twierdzenie \ref{ma2:thm:inverse-perturbation}
w algebrze Banacha $\mathcal L(E)$ z elementem $I_E$
i zaburzeniem $-tK$, którego norma jest mniejsza od $1$.
Otrzymujemy odwracalność $I_E-tK$ i jego odwrotność
w postaci szeregu $\sum_{n=0}^\infty t^nK^n$ w normie
operatorowej. Zastosowanie operatora do ustalonego $b$
jest ciągłe, zatem jedyne rozwiązanie w całym $E$ wynosi
\[
 u(t)=\sum_{n=0}^\infty t^n K^n b.
\]
Szereg jest zbieżny w $E$, gdyż norma $n$-tego
składnika jest ograniczona przez
$(|t|\|K\|)^n\|b\|$ i $E$ jest zupełna.
To rozwinięcie jest zgodne z przykładem
\atlasref{FA-15::ex:implicit-operator15}{FA-15: rodzina równań operatorowych}.
\end{przyklad}
